\section{Experimental Details}
\label{app:experimental-details}

This appendix records the data, training, and evaluation settings used for the
experiments in the main text. Unless a table or paragraph states otherwise, all
reported \opsd{} runs use the defaults in Table~\ref{tab:opsd-training-hparams}
and all evaluations use the generation and grading protocol in
Table~\ref{tab:eval-hparams}.

\subsection{Data}
\label{app:data}

\paragraph{OpenThoughts math.}
We use a cleaned version of \texttt{Ashkchamp/Openthoughts\_math\_filtered\_30K}.
The cleaning script removes system turns, remaps thought delimiters to
\texttt{<think>} and \texttt{</think>}, removes explicit solution delimiters,
and appends the instruction to return the final answer in \texttt{\textbackslash
boxed\{\}}. We use a 15K subset of the data for training. The privileged teacher context is taken from the \texttt{solution} column.

\paragraph{Countdown.}
For Countdown, we use \texttt{jasonrqh/Countdown-CoT-20k}. We select a 15K subset for training and reserve an additional
500 examples as a held-out evaluation set.

\begin{table}[!htbp]
    \centering
    \small
    \caption{Training and evaluation data used in the experiments.}
    \label{tab:data-summary}
    \begin{tabular}{ll>{\raggedright\arraybackslash}p{0.58\linewidth}}
        \toprule
        Data & Rows & Use and fields \\
        \midrule
        OpenThoughts cleaned & 29,439 &
        Source pool after cleaning; columns include \texttt{problem},
        \texttt{solution}, and \texttt{Answer}. \\
        OpenThoughts 15k & 15,000 &
        Main OpenThoughts \opsd{} training subset; prompt =
        \texttt{problem}, privileged context = \texttt{solution}. \\
        Countdown train & 15,000 &
        Countdown \opsd{} training; prompt =
        \texttt{datapoint\_input\_text}, privileged context =
        \texttt{response\_suffix}. \\
        Countdown eval & 500 &
        Held-out Countdown evaluation. \\
        AIME24/AIME25/HMMT25 & 30 each &
        Math competition evaluations. \\
        \bottomrule
    \end{tabular}
\end{table}

\subsection{Existing Asset Licenses}
\label{app:asset-licenses}

Table~\ref{tab:asset-licenses} summarizes the existing datasets, model
checkpoints, and software assets used in our experiments. We use these assets
for training, evaluation, or implementation, and do not redistribute
third-party datasets or model checkpoints.

\begin{table}[!htbp]
    \centering
    \scriptsize
    \caption{Existing assets used in this work.}
    \label{tab:asset-licenses}
    \setlength{\tabcolsep}{3pt}
    \renewcommand{\arraystretch}{1.15}
    \begin{tabular}{>{\raggedright\arraybackslash}p{0.23\linewidth}
                    >{\raggedright\arraybackslash}p{0.28\linewidth}
                    >{\raggedright\arraybackslash}p{0.18\linewidth}
                    >{\raggedright\arraybackslash}p{0.23\linewidth}}
        \toprule
        Asset & Source & License / terms & Use \\
        \midrule
        \texttt{Ashkchamp/Openthoughts\_math\_filtered\_30K} &
        \url{https://huggingface.co/datasets/Ashkchamp/Openthoughts_math_filtered_30K} &
        No explicit license metadata listed on the dataset page at time of access &
        Source pool for the OpenThoughts math training subset; not redistributed. \\
        \texttt{open-thoughts/OpenThoughts-114k} &
        \url{https://huggingface.co/datasets/open-thoughts/OpenThoughts-114k} &
        Apache-2.0 &
        Upstream OpenThoughts dataset related to the math training data; not redistributed. \\
        \texttt{jasonrqh/Countdown-CoT-20k} &
        \url{https://huggingface.co/datasets/jasonrqh/Countdown-CoT-20k} &
        MIT &
        Countdown training and held-out in-domain evaluation data. \\
        Qwen3 checkpoints &
        \url{https://huggingface.co/Qwen} &
        Apache-2.0 &
        Base instruction and thinking models used for training and evaluation. \\
        OLMo-3 checkpoints &
        \url{https://huggingface.co/allenai/Olmo-3-7B-Think} &
        Apache-2.0; Ai2 Responsible Use Guidelines &
        Base instruction and thinking models used for training and evaluation. \\
        \texttt{idanshen/Self-Distillation} &
        \url{https://github.com/idanshen/Self-Distillation} &
        No explicit software license file found at time of access &
        Reference/adapted implementation code for self-distillation; not redistributed. \\
        \texttt{vLLM} &
        \url{https://github.com/vllm-project/vllm} &
        Apache-2.0 &
        Inference engine used for generation and evaluation. \\
        AIME and HMMT problems &
        Public competition materials &
        Public competition materials; not redistributed &
        Out-of-domain math evaluation benchmarks. \\
        \bottomrule
    \end{tabular}
\end{table}

\subsection{Prompt Templates and Privileged Context Examples}
\label{app:prompt-templates}

\newcommand{\promptvar}[1]{\texttt{\textcolor{RoyalBlue!70!black}{\{#1\}}}}
\newcommand{\promptcard}[3]{%
    \begin{center}
    \setlength{\fboxsep}{7pt}%
    \fcolorbox{#1!55!black}{#1!5}{%
        \begin{minipage}{0.92\linewidth}
        \raggedright
        \textbf{\textcolor{#1!55!black}{#2}}\par\vspace{4pt}
        \small #3
        \end{minipage}%
    }%
    \end{center}
}

We fill the following templates before applying each model's chat template.
For standard \opsd{} training, the student receives only the problem.
\promptcard{RoyalBlue}{Student prompt template}{%
{\ttfamily
\promptvar{problem}
}
}

The privileged teacher receives the same problem plus an example response.
\promptcard{OliveGreen}{Privileged-teacher prompt template}{%
{\ttfamily
\promptvar{problem}\par\medskip
This is an example for a response to the question:\par\medskip
\promptvar{Answer}\par\medskip
Now answer with a response of your own, including the thinking process.
}
}
Here \promptvar{Answer} denotes the privileged response used as teacher
context. In dense gold-demonstration runs, this field contains a full reference
solution; in sparse final-answer-only runs, it contains only the boxed final
answer.

For the conciseness-control experiment in
Figure~\ref{fig:openthoughts_qwen8b_concise_budget}, the teacher is prompted
without gold context.
\promptcard{BurntOrange}{Conciseness-control teacher prompt}{%
{\ttfamily
Solve the following math problem concisely and correctly. Be direct -- avoid
unnecessary elaboration, redundant steps, or restating the problem. Focus only
on the key reasoning steps needed to reach the answer.\par\medskip
\promptvar{problem}
}
}

\clearpage
\paragraph{Sparse and dense context example.}
The dense and sparse conditions use the same privileged-teacher wrapper above;
they differ only in what is inserted into the \promptvar{Answer} slot. The
following filled examples use the same problem.

\promptcard{OliveGreen}{Dense teacher prompt, filled}{%
{\ttfamily
Given real numbers \(a, b, c\) and a positive number \(\lambda\) such that the
polynomial \(f(x) = x^3 + a x^2 + b x + c\) has three real roots \(x_1, x_2,
x_3\), and the conditions \(x_2 - x_1 = \lambda\) and
\(x_3 > \frac{1}{2}(x_1 + x_2)\) are satisfied, find the maximum value of
\[
\frac{2 a^3 + 27 c - 9 a b}{\lambda^3}.
\]
\par\medskip
This is an example for a response to the question:\par\medskip
We begin by analyzing the function \(f(x) = x^3 + a x^2 + b x + c\), which has
three real roots \(x_1, x_2, x_3\). We are given the following conditions:
\(x_2 - x_1 = \lambda\) and
\(x_3 > \frac{1}{2}(x_1 + x_2)\). We aim to find the maximum value of
\[
\frac{2a^3 + 27c - 9ab}{\lambda^3}.
\]
\par\medskip
Transform the polynomial to remove the quadratic term.
Substitute \(x = y - \frac{a}{3}\) into \(f(x)\):
\[
\begin{array}{rcl}
F(y) & = & f\left(y - \frac{a}{3}\right) \\
     & = & \left(y - \frac{a}{3}\right)^3
           + a \left(y - \frac{a}{3}\right)^2
           + b \left(y - \frac{a}{3}\right) + c \\
     & = & y^3 - \left(\frac{a^2}{3} - b\right)y
           + \frac{1}{27}(2a^3 + 27c - 9ab).
\end{array}
\]
\par\medskip
Identify the new roots of \(F(y)\).
Let the roots of \(F(y)\) be \(y_1, y_2, y_3\). We know
\(y_i = x_i + \frac{a}{3}\). Using Vieta's formulas,
\[
y_1 + y_2 + y_3 = 0,\qquad
y_1 y_2 y_3 = -\frac{1}{27}(2a^3 + 27c - 9ab).
\]
\par\medskip
\textit{[Middle of the gold demonstration omitted for clarity of the example.]}
\par\medskip
Conclusion:
\[
\boxed{\frac{3\sqrt{3}}{2}}
\]
\par\medskip
Now answer with a response of your own, including the thinking process.
}
}

\promptcard{OliveGreen}{Sparse final-answer-only teacher prompt, filled}{%
{\ttfamily
Given real numbers \(a, b, c\) and a positive number \(\lambda\) such that the
polynomial \(f(x) = x^3 + a x^2 + b x + c\) has three real roots \(x_1, x_2,
x_3\), and the conditions \(x_2 - x_1 = \lambda\) and
\(x_3 > \frac{1}{2}(x_1 + x_2)\) are satisfied, find the maximum value of
\[
\frac{2 a^3 + 27 c - 9 a b}{\lambda^3}.
\]
\par\medskip
This is an example for a response to the question:\par
\[
\boxed{\frac{3\sqrt{3}}{2}}
\]
\par\medskip
Now answer with a response of your own, including the thinking process.
}
}

\subsection{OPSD Training}
\label{app:opsd-training}

For each training example, the student is prompted with the task input alone.
The teacher is initialized from the same base checkpoint, but receives the same
task input plus privileged supervision. The trainer samples completions
on-policy and minimizes token-level divergence between teacher and student
distributions on those sampled completion tokens.

\begin{table}[!htbp]
    \centering
    \small
    \caption{Default \opsd{} training hyperparameters.}
    \label{tab:opsd-training-hparams}
    \begin{tabular}{ll}
        \toprule
        Hyperparameter & Value \\
        \midrule
        Objective & Generalized KL/JSD distillation on sampled completion tokens \\
        Distillation mixture $\alpha$ & 0.5 (JSD) \\
        Teacher synchronization & Disabled for reported regular runs \\
        Completions per prompt & 1 \\
        Training sampling temperature / top-$p$  & 1.0 / 1.0  \\
        Epochs & 1 \\
        Effective batch size & 64 \\
        Per-device batch size & 1 \\
        Gradient accumulation & 8 \\
        Optimizer learning-rate schedule & Cosine decay \\
        Warmup ratio & 0.1 \\
        Max gradient norm & 1.0 \\
        Weight decay & 0.0 \\
        Precision & bf16 \\
        Distributed training & FSDP full-shard auto-wrap \\
        Gradient checkpointing & Enabled \\
        Seed & 42 \\
        Max prompt length & 25,000 tokens \\
        Max completion length & 4,096 tokens unless noted otherwise \\
        \bottomrule
    \end{tabular}
\end{table}

\begin{table}[!htbp]
    \centering
    \small
    \caption{Model-specific \opsd{} training settings and deviations from
    Table~\ref{tab:opsd-training-hparams}.}
    \label{tab:model-training-hparams}
    \resizebox{\textwidth}{!}{%
    \begin{tabular}{llllll}
        \toprule
        Model(s) & Training data & LR & Max completion & vLLM train gen. & Hardware \\
        \midrule
        Qwen3-1.7B, Qwen3-4B, Qwen3-4B-Thinking-2507,
        Qwen3-4B-Instruct-2507 &
        OpenThoughts 15k & $5{\times}10^{-6}$ & 4,096 & Yes &
        8 H100, 80GB/GPU \\
        Qwen3-8B &
        OpenThoughts 15k & $2{\times}10^{-6}$ & 4,096 & Yes &
        8 H100, 80GB/GPU \\
        OLMo-3-7B-Instruct, OLMo-3-7B-Think &
        OpenThoughts 15k & $5{\times}10^{-6}$ & 4,096 & Yes &
        8 H100, 80GB/GPU \\
        Qwen3-4B-Thinking-2507, Qwen3-4B-Instruct-2507,
        OLMo-3-7B-Instruct, OLMo-3-7B-Think &
        Countdown 15k & $5{\times}10^{-6}$ & 4,096 & Yes &
        8 H100, 80GB/GPU \\
        \bottomrule
    \end{tabular}%
    }
\end{table}

\subsection{Evaluation}
\label{app:evaluation}

For AIME24, AIME25, and HMMT25, we generate 16 samples per problem on 30
problems per benchmark. Generation is sharded over 8 jobs and uses a maximum
generation length of 38,912 tokens. The merged generation file therefore
contains 480 rows for each AIME/HMMT benchmark. Countdown uses the same
16-sample evaluation protocol on the 500-example held-out split when reported.

\begin{table}[!htbp]
    \centering
    \small
    \caption{Evaluation generation hyperparameters.}
    \label{tab:eval-hparams}
    \begin{tabular}{lll}
        \toprule
        Model/eval group & Temperature & Top-$p$ \\
        \midrule
        Qwen3 thinking-style models & 0.6 & 0.95 \\
        Qwen3-4B-Instruct-2507 & 0.7 & 0.8 \\
        OLMo-3-7B-Instruct/Think & 0.6 & 0.95 \\
        \bottomrule
    \end{tabular}
\end{table}

\begin{table}[!htbp]
    \centering
    \small
    \caption{Shared evaluation settings.}
    \label{tab:eval-shared-hparams}
    \begin{tabular}{ll}
        \toprule
        Setting & Value \\
        \midrule
        Samples per problem & 16 \\
        Shards & 8 \\
        Maximum generation length & 38,912 tokens \\
        Eval engine & vLLM \\
        Eval precision & bf16 \\
        vLLM GPU memory utilization & 0.9 \\
        Eval top-$k$ & -1 \\
        Prompt format & Model chat template with generation prompt \\
        Metric file & \texttt{metrics.json} written next to merged JSONL \\
        \bottomrule
    \end{tabular}
\end{table}

\paragraph{Metrics.}
The reported avg@16 \texttt{accuracy} is the mean correctness over generated
samples: for each problem, we average correctness across its 16 sampled rollouts,
then average across problems. This is equivalent to empirical single-sample
correctness under the evaluation sampling distribution, but is distinct from
unbiased pass@16.
For $k > 1$, pass@$k$ is computed
with the unbiased estimator. For each problem with $n$ samples and $c$ correct
samples, the contribution is 1 if $n-c < k$ and otherwise
\[
1 - \prod_{i=0}^{k-1} \frac{n-c-i}{n-i}.
\]
We report pass@$k$ only for $k \leq n$.

\subsection{Deliberation Marker Analysis}
\label{app:deliberation-markers}

To test whether \opsd{} changes explicit deliberation language in model
reasoning, we compare paired base and \opsd{} rollouts for the same model,
benchmark, problem, and sample (see Table~\ref{tab:deliberation_markers_full}). The analysis uses the five thinking models
in the OpenThoughts 15k comparison on AIME24, AIME25, and HMMT25. With 16
samples for each of 30 problems on each benchmark, this gives 7,200 paired
rollouts.

For each response, we count occurrences of a fixed, case-insensitive lexicon of
deliberation markers. The marker families are verification markers, with
examples such as \textit{check}, \textit{verify}, \textit{double check}, and
\textit{make sure}; backtracking markers, with examples such as \textit{wait},
\textit{actually}, \textit{mistake}, \textit{wrong}, and \textit{another way};
and hedging markers, with examples such as \textit{maybe}, \textit{probably},
\textit{might}, and \textit{seems}. We report raw marker counts and marker
counts per 1,000 response tokens. Response-token counts are computed with the
cached model tokenizer.

\begin{table*}[t]
\centering
\caption{\textbf{Full deliberation-marker counts for paired base and \opsd{} rollouts.}
Raw columns report average marker counts per response. Normalized columns report occurrences per 1,000 generated tokens, using response-token counts computed with the cached model tokenizer. Deltas are \opsd{} minus base, with confidence intervals computed by clustered bootstrap over model--benchmark--problem clusters.}
\label{tab:deliberation_markers_full}
\small
\setlength{\tabcolsep}{4pt}
\renewcommand{\arraystretch}{1.12}
\begin{tabular*}{\textwidth}{@{\extracolsep{\fill}}lcccccc@{}}
\toprule
\textbf{Marker family}
& \textbf{Base raw}
& \textbf{\opsd{} raw}
& \textbf{Raw $\Delta$}
& \textbf{Base /1k}
& \textbf{\opsd{} /1k}
& \textbf{$\Delta$ /1k} \\
\midrule
Verification
& 26.1
& 17.3
& $-8.8$ [$-9.5$, $-8.1$]
& 1.63
& 1.35
& $-0.28$ [$-0.32$, $-0.25$] \\
Backtracking
& 124.2
& 100.6
& $-23.5$ [$-29.0$, $-19.1$]
& 6.45
& 6.29
& $-0.16$ [$-0.32$, $-0.02$] \\
Hedging
& 72.2
& 55.2
& $-17.0$ [$-18.5$, $-15.5$]
& 3.79
& 3.62
& $-0.17$ [$-0.25$, $-0.10$] \\
\bottomrule
\end{tabular*}
\end{table*}


Deltas are \opsd{} minus base. Confidence intervals are computed by clustered
bootstrap over model--benchmark--problem clusters, so the 16 samples from the
same problem are not treated as fully independent. The mean response length in
this analysis is 19,395 tokens for base rollouts and 15,391 tokens after
\opsd{}. Since the counts are lexical proxies, we interpret them as evidence
about explicit deliberation markers in the generated traces, not as direct
measurements of latent uncertainty or confidence.
